% frontmatter/polish/actividades.tex -- "Zadania": las de
% frontmatter/actividades.tex, en polaco. Tiene que seguir cabiendo en una
% página (tools/check_pages.py comprueba que el día 1 está en la página
% 5).
\section*{Zadania}

Każde zadanie pojawia się w dniu, w którym jest potrzebne po raz
pierwszy, a potem co jakiś czas wraca. Polecenia na każdej stronie, małą
szarą czcionką, czyta dorosły.

{\small
\begin{multicols}{2}
\raggedright
\subsection*{Od jesieni}
\begin{itemize}[leftmargin=5mm, itemsep=1pt, topsep=0pt]
\item \textbf{\lblTraza} (w poniedziałki, kiedy przychodzi nowa liczba):
  pisać liczbę palcem po kropkach, potem ołówkiem, a potem samodzielnie
  na linii. Kropki to kształt prawdziwej liczby, tej samej, którą widać
  w całym zeszycie.
\item \textbf{\lblColorea}: pokolorować dokładnie tyle, ile jest
  napisane, ani jednej więcej -- reszta zostaje biała.
\item \textbf{\lblRodea}: spośród kilku grup rzeczy zakreślić tę, w
  której jest tyle, ile jest napisane. Do 6 grupy są ułożone jak oczka na
  kostce, a od 7 do 10 jak w ramce dziesiątkowej, żeby było je widać od
  razu; pusta grupa to 0.
\item \textbf{\lblCuenta}: policzyć, wskazując każdą rzecz raz, i
  zakreślić liczbę.
\item \textbf{\lblBusca}: zakreślić liczbę za każdym razem, kiedy się
  pojawia, między innymi liczbami i kształtami.
\item \textbf{\lblUne}: policzyć każdą grupę i poprowadzić linię od jej
  kropki do kropki jej liczby.
\item \textbf{\lblCompleta}: mówić na głos liczby z pociągu, po kolei, i
  wpisać brakujące do pustych wagonów.
\item \textbf{\lblDibuja}: narysować dokładnie to, co jest napisane, ani
  więcej, ani mniej.
\item \textbf{\lblRepasa} (w piątki): czego nauczyliśmy się w tym
  tygodniu, do zaznaczenia, i rysunek. Co dziesięć stron jest mały napis,
  żeby to uczcić.
\end{itemize}

\subsection*{Od zimy}
\begin{itemize}[leftmargin=5mm, itemsep=1pt, topsep=0pt]
\item \textbf{\lblDecena} (od 11 do 20): pełna taca, czyli 10, i druga
  obok; liczyć dalej od 10.
\item \textbf{\lblCompara}: zakreślić grupę, w której jest więcej albo
  mniej, albo dwie, w których jest tyle samo; albo największą liczbę,
  albo najmniejszą.
\item \textbf{\lblVecinos}: wpisać liczbę, która jest przed, i tę, która
  jest po, w domkach po obu stronach.
\item \textbf{\lblRecta}: mówić liczby z osi, od lewej do prawej, i
  wpisać brakujące.
\item \textbf{\lblOrdena}: wpisać liczby po kolei, od najmniejszej do
  największej (albo odwrotnie).
\item \textbf{\lblJunta}: połączyć dwie grupy i policzyć, ile jest
  razem: dodawanie, jeszcze bez znaków + i =.
\end{itemize}

\subsection*{Od wiosny}
\begin{itemize}[leftmargin=5mm, itemsep=1pt, topsep=0pt]
\item \textbf{\lblSuma} i \textbf{\lblResta}: to samo, już ze znakami;
  przy odejmowaniu skreśla się to, co zabieramy.
\item \textbf{\lblParte}: kreska dzieli grupę na dwie części; policzyć
  to, co jest po drugiej stronie.
\item \textbf{\lblDiez}: dorysować kropki, których brakuje, żeby
  wypełnić ramkę dziesiątkową.
\item \textbf{\lblProblema}: czyta je dorosły; można je narysować, a
  potem zapisać działanie.
\item \textbf{\lblBloques}: każdy słupek to dziesiątka (dziesięć
  kostek); policzyć słupki i pojedyncze kostki.
\item \textbf{\lblCompleta}, \textbf{\lblRecta} i \textbf{\lblCuenta}
  także dwójkami, piątkami i dziesiątkami: koła rowerów, palce u rąk,
  dziesiątki.
\end{itemize}

\subsection*{Od lata}
\begin{itemize}[leftmargin=5mm, itemsep=1pt, topsep=0pt]
\item \textbf{\lblTabla}: każdy rząd to dziesiątka; mówić liczby rząd po
  rzędzie i wpisać brakujące.
\item \textbf{\lblDinero}: policzyć najpierw banknoty, a potem monety, i
  wpisać, ile jest euro (Lucía mieszka w Hiszpanii).
\item \textbf{\lblHora}: krótka wskazówka pokazuje godzinę; kiedy długa
  jest na samej górze, jest pełna godzina, a kiedy na samym dole -- wpół
  do następnej (wpół do czwartej to trzecia trzydzieści).
\item \textbf{\lblMide}: policzyć kostki pod ołówkiem, od gumki do
  czubka; przy dwóch ołówkach zakreślić dłuższy.
\end{itemize}
\end{multicols}}
\newpage
